% Options for packages loaded elsewhere
% Options for packages loaded elsewhere
\PassOptionsToPackage{unicode}{hyperref}
\PassOptionsToPackage{hyphens}{url}
\PassOptionsToPackage{dvipsnames,svgnames,x11names}{xcolor}
%
\documentclass[
  11pt,
  letterpaper,
  DIV=11,
  numbers=noendperiod]{scrartcl}
\usepackage{xcolor}
\usepackage[margin=1in]{geometry}
\usepackage{amsmath,amssymb}
\setcounter{secnumdepth}{5}
\usepackage{iftex}
\ifPDFTeX
  \usepackage[T1]{fontenc}
  \usepackage[utf8]{inputenc}
  \usepackage{textcomp} % provide euro and other symbols
\else % if luatex or xetex
  \usepackage{unicode-math} % this also loads fontspec
  \defaultfontfeatures{Scale=MatchLowercase}
  \defaultfontfeatures[\rmfamily]{Ligatures=TeX,Scale=1}
\fi
\usepackage{lmodern}
\ifPDFTeX\else
  % xetex/luatex font selection
\fi
% Use upquote if available, for straight quotes in verbatim environments
\IfFileExists{upquote.sty}{\usepackage{upquote}}{}
\IfFileExists{microtype.sty}{% use microtype if available
  \usepackage[]{microtype}
  \UseMicrotypeSet[protrusion]{basicmath} % disable protrusion for tt fonts
}{}
\makeatletter
\@ifundefined{KOMAClassName}{% if non-KOMA class
  \IfFileExists{parskip.sty}{%
    \usepackage{parskip}
  }{% else
    \setlength{\parindent}{0pt}
    \setlength{\parskip}{6pt plus 2pt minus 1pt}}
}{% if KOMA class
  \KOMAoptions{parskip=half}}
\makeatother
% Make \paragraph and \subparagraph free-standing
\makeatletter
\ifx\paragraph\undefined\else
  \let\oldparagraph\paragraph
  \renewcommand{\paragraph}{
    \@ifstar
      \xxxParagraphStar
      \xxxParagraphNoStar
  }
  \newcommand{\xxxParagraphStar}[1]{\oldparagraph*{#1}\mbox{}}
  \newcommand{\xxxParagraphNoStar}[1]{\oldparagraph{#1}\mbox{}}
\fi
\ifx\subparagraph\undefined\else
  \let\oldsubparagraph\subparagraph
  \renewcommand{\subparagraph}{
    \@ifstar
      \xxxSubParagraphStar
      \xxxSubParagraphNoStar
  }
  \newcommand{\xxxSubParagraphStar}[1]{\oldsubparagraph*{#1}\mbox{}}
  \newcommand{\xxxSubParagraphNoStar}[1]{\oldsubparagraph{#1}\mbox{}}
\fi
\makeatother


\usepackage{longtable,booktabs,array}
\newcounter{none} % for unnumbered tables
\usepackage{calc} % for calculating minipage widths
% Correct order of tables after \paragraph or \subparagraph
\usepackage{etoolbox}
\makeatletter
\patchcmd\longtable{\par}{\if@noskipsec\mbox{}\fi\par}{}{}
\makeatother
% Allow footnotes in longtable head/foot
\IfFileExists{footnotehyper.sty}{\usepackage{footnotehyper}}{\usepackage{footnote}}
\makesavenoteenv{longtable}
\usepackage{graphicx}
\makeatletter
\newsavebox\pandoc@box
\newcommand*\pandocbounded[1]{% scales image to fit in text height/width
  \sbox\pandoc@box{#1}%
  \Gscale@div\@tempa{\textheight}{\dimexpr\ht\pandoc@box+\dp\pandoc@box\relax}%
  \Gscale@div\@tempb{\linewidth}{\wd\pandoc@box}%
  \ifdim\@tempb\p@<\@tempa\p@\let\@tempa\@tempb\fi% select the smaller of both
  \ifdim\@tempa\p@<\p@\scalebox{\@tempa}{\usebox\pandoc@box}%
  \else\usebox{\pandoc@box}%
  \fi%
}
% Set default figure placement to htbp
\def\fps@figure{htbp}
\makeatother





\setlength{\emergencystretch}{3em} % prevent overfull lines

\providecommand{\tightlist}{%
  \setlength{\itemsep}{0pt}\setlength{\parskip}{0pt}}



 


\KOMAoption{captions}{tableheading}
\makeatletter
\@ifpackageloaded{caption}{}{\usepackage{caption}}
\AtBeginDocument{%
\ifdefined\contentsname
  \renewcommand*\contentsname{Table of contents}
\else
  \newcommand\contentsname{Table of contents}
\fi
\ifdefined\listfigurename
  \renewcommand*\listfigurename{List of Figures}
\else
  \newcommand\listfigurename{List of Figures}
\fi
\ifdefined\listtablename
  \renewcommand*\listtablename{List of Tables}
\else
  \newcommand\listtablename{List of Tables}
\fi
\ifdefined\figurename
  \renewcommand*\figurename{Figure}
\else
  \newcommand\figurename{Figure}
\fi
\ifdefined\tablename
  \renewcommand*\tablename{Table}
\else
  \newcommand\tablename{Table}
\fi
}
\@ifpackageloaded{float}{}{\usepackage{float}}
\floatstyle{ruled}
\@ifundefined{c@chapter}{\newfloat{codelisting}{h}{lop}}{\newfloat{codelisting}{h}{lop}[chapter]}
\floatname{codelisting}{Listing}
\newcommand*\listoflistings{\listof{codelisting}{List of Listings}}
\makeatother
\makeatletter
\makeatother
\makeatletter
\@ifpackageloaded{caption}{}{\usepackage{caption}}
\@ifpackageloaded{subcaption}{}{\usepackage{subcaption}}
\makeatother
\usepackage{bookmark}
\IfFileExists{xurl.sty}{\usepackage{xurl}}{} % add URL line breaks if available
\urlstyle{same}
\hypersetup{
  pdftitle={U.S. Support to OCHA Country-Based Pooled Funds},
  pdfauthor={Spencer},
  colorlinks=true,
  linkcolor={blue},
  filecolor={Maroon},
  citecolor={Blue},
  urlcolor={Blue},
  pdfcreator={LaTeX via pandoc}}


\title{U.S. Support to OCHA Country-Based Pooled Funds}
\usepackage{etoolbox}
\makeatletter
\providecommand{\subtitle}[1]{% add subtitle to \maketitle
  \apptocmd{\@title}{\par {\large #1 \par}}{}{}
}
\makeatother
\subtitle{Tracking Publicly Announced FY2026 U.S. Contributions Through
the UN OCHA CBPF Data Platform}
\author{Spencer}
\date{2026-07-03}
\begin{document}
\maketitle

\renewcommand*\contentsname{Table of contents}
{
\setcounter{tocdepth}{3}
\tableofcontents
}

\section{Executive Summary}\label{executive-summary}

\begin{quote}
\textbf{Draft}
\end{quote}

This report examines publicly announced U.S. Government (USG) support to
the United Nations Office for the Coordination of Humanitarian Affairs
(OCHA) Country-Based Pooled Funds (CBPFs) during Fiscal Year (FY) 2026.
Using official OCHA API data, the report traces announced U.S.
contributions to individual pooled funds and the humanitarian projects
associated with those funds.

The analysis compares the December 29, 2025 Memorandum of Understanding
(MOU) and the May 14, 2026 supplemental funding announcement with live
administrative data published through the OCHA CBPF API.

Key findings include:

\begin{itemize}
\tightlist
\item
  Total U.S. pledged funding
\item
  Total U.S. paid funding
\item
  Number of pooled funds supported
\item
  Number of associated humanitarian projects
\item
  Major implementing organizations
\item
  Geographic distribution of funding
\end{itemize}

\begin{center}\rule{0.5\linewidth}{0.5pt}\end{center}

\section{1. Background}\label{background}

\subsection{Country-Based Pooled
Funds}\label{country-based-pooled-funds}

Describe what CBPFs are.

Describe how OCHA administers pooled funding.

Describe why pooled funds are important humanitarian financing
mechanisms.

\subsection{U.S.--OCHA Partnership}\label{u.s.ocha-partnership}

Describe the December 29, 2025 MOU.

Describe the May 14, 2026 supplemental announcement.

Explain that these announcements form the basis for this analysis.

\begin{center}\rule{0.5\linewidth}{0.5pt}\end{center}

\section{2. Research Questions}\label{research-questions}

This report seeks to answer the following questions:

\begin{enumerate}
\def\labelenumi{\arabic{enumi}.}
\tightlist
\item
  Which CBPFs received FY2026 U.S. contributions?
\item
  Which countries announced by the United States appear in official OCHA
  contribution data?
\item
  How much funding has been pledged and paid?
\item
  Which humanitarian projects are associated with those pooled funds?
\item
  Which organizations are implementing those projects?
\end{enumerate}

\begin{center}\rule{0.5\linewidth}{0.5pt}\end{center}

\section{3. Data Sources}\label{data-sources}

\subsection{Official Sources}\label{official-sources}

\begin{itemize}
\tightlist
\item
  OCHA CBPF Contribution API
\item
  OCHA CBPF Project Summary API
\item
  December 29, 2025 MOU announcement
\item
  May 14, 2026 supplemental funding announcement
\end{itemize}

\subsection{Methodology}\label{methodology}

This report filters contribution data using:

\begin{itemize}
\tightlist
\item
  Donor = United States
\item
  Fiscal Year = 2026
\end{itemize}

Projects are linked using \textbf{PooledFundId}, the common identifier
shared across OCHA endpoints.

\begin{center}\rule{0.5\linewidth}{0.5pt}\end{center}

\section{4. Overview of U.S.
Commitments}\label{overview-of-u.s.-commitments}

\subsection{December 2025 MOU}\label{december-2025-mou}

\emph{Insert narrative}

\subsubsection{Countries}\label{countries}

\emph{Insert table}

\subsection{May 2026 Supplemental
Funding}\label{may-2026-supplemental-funding}

\emph{Insert narrative}

\subsubsection{Countries}\label{countries-1}

\emph{Insert table}

\subsection{Comparison}\label{comparison}

Discuss countries retained and countries added.

\begin{center}\rule{0.5\linewidth}{0.5pt}\end{center}

\section{5. Crosswalk Between Public Announcements and OCHA
Data}\label{crosswalk-between-public-announcements-and-ocha-data}

This chapter compares official announcements with the live OCHA CBPF
API.

Insert master crosswalk table.

Questions addressed:

\begin{itemize}
\tightlist
\item
  Which countries appear in both announcements?
\item
  Which countries appear in OCHA data?
\item
  Which countries have associated projects?
\end{itemize}

\begin{center}\rule{0.5\linewidth}{0.5pt}\end{center}

\section{6. U.S. Contributions by Pooled
Fund}\label{u.s.-contributions-by-pooled-fund}

Insert:

\begin{itemize}
\tightlist
\item
  Bar chart of pledged funding
\item
  Bar chart of paid funding
\end{itemize}

Discuss:

\begin{itemize}
\tightlist
\item
  Largest recipients
\item
  Regional pooled funds
\item
  Differences between pledged and paid amounts
\end{itemize}

\begin{center}\rule{0.5\linewidth}{0.5pt}\end{center}

\section{7. Humanitarian Projects
Supported}\label{humanitarian-projects-supported}

Summarize:

\begin{itemize}
\tightlist
\item
  Number of projects
\item
  Total project budget
\item
  Average project budget
\end{itemize}

Insert:

\begin{itemize}
\tightlist
\item
  Table of projects by pooled fund
\item
  Table of project counts
\end{itemize}

\begin{center}\rule{0.5\linewidth}{0.5pt}\end{center}

\section{8. Implementing
Organizations}\label{implementing-organizations}

Summarize:

\begin{itemize}
\tightlist
\item
  Largest implementing partners
\item
  Organization types
\item
  Number of projects by partner
\end{itemize}

Insert charts.

\begin{center}\rule{0.5\linewidth}{0.5pt}\end{center}

\section{9. Geographic Distribution}\label{geographic-distribution}

Insert:

\begin{itemize}
\tightlist
\item
  Country summary table
\item
  Map (future)
\end{itemize}

Discuss geographic coverage.

\begin{center}\rule{0.5\linewidth}{0.5pt}\end{center}

\section{10. Key Findings}\label{key-findings}

Summarize major observations.

Example topics:

\begin{itemize}
\tightlist
\item
  Coverage of the original MOU
\item
  Expansion under the May announcement
\item
  Distribution across pooled funds
\item
  Largest humanitarian operations
\item
  Relationship between contributions and approved projects
\end{itemize}

\begin{center}\rule{0.5\linewidth}{0.5pt}\end{center}

\section{11. Limitations}\label{limitations}

Important considerations:

\begin{itemize}
\tightlist
\item
  CBPFs are pooled financing mechanisms.
\item
  Individual projects cannot be interpreted as receiving direct U.S.
  funding.
\item
  Project approval dates and contribution dates may differ.
\item
  API data represents a snapshot at the time of analysis.
\end{itemize}

\begin{center}\rule{0.5\linewidth}{0.5pt}\end{center}

\section{12. Conclusion}\label{conclusion}

Summarize the overall findings.

Discuss opportunities for future analysis.

\begin{center}\rule{0.5\linewidth}{0.5pt}\end{center}

\section{Appendix A --- Methodology}\label{appendix-a-methodology}

Describe:

\begin{itemize}
\tightlist
\item
  API endpoints
\item
  Filters
\item
  Country normalization
\item
  Matching using PooledFundId
\end{itemize}

\begin{center}\rule{0.5\linewidth}{0.5pt}\end{center}

\section{Appendix B --- Code}\label{appendix-b-code}

Include the Python workflow used to reproduce the analysis.

\begin{center}\rule{0.5\linewidth}{0.5pt}\end{center}

\section{Appendix C --- Data
Dictionary}\label{appendix-c-data-dictionary}

Describe major variables used.

Examples:

{\def\LTcaptype{none} % do not increment counter
\begin{longtable}[]{@{}ll@{}}
\toprule\noalign{}
Variable & Description \\
\midrule\noalign{}
\endhead
\bottomrule\noalign{}
\endlastfoot
PooledFundId & Unique pooled fund identifier \\
PooledFundName & OCHA pooled fund name \\
DonorName & Contributing donor \\
FiscalYear & Contribution fiscal year \\
PledgeAmt & Amount pledged \\
PaidAmt & Amount paid \\
AllocationYear & Project allocation year \\
Budget & Approved project budget \\
\end{longtable}
}

\begin{center}\rule{0.5\linewidth}{0.5pt}\end{center}

\section{References}\label{references}

United Nations Office for the Coordination of Humanitarian Affairs.

Country-Based Pooled Funds API.

United States Department of State.

United Nations News.

Additional references.




\end{document}
